\begin{example}[Truth without knowledge]
Let \(w\sim_i v\), let \(p\) be true at \(w\), and let \(p\) be false at
\(v\).  Then \(p\) is true at \(w\) but \(K_ip\) is false there.  The structure
does not refute factivity \(K_ip\to p\): equivalence makes \(w\sim_i w\), so
knowledge still implies truth at the actual world.
\end{example}

\begin{warning}[Changing the accessibility condition]
If an epistemic relation is not reflexive, \(K_iA\to A\) can fail.  Such a
modality may represent belief or information rather than factive knowledge, but
the change must be made in both terminology and frame class.
\end{warning}

